% GENERATED FILE. Edit canonical lessons, then run npm run generate:books.
% canonical-source-hash: fnv1a64:3bb639b2a53506e8
% canonical-lessons: LA-C304-effringo, LA-C304-claustrum, LA-C304-rapio, LA-C304-deprehendo, LA-C304-comprehendo

\chapter{Fūr comprehēnsus}
\label{ch:la-fur-comprehensus}
\hlchaptermodality{304}

\begin{chapteropening}
\textbf{By the end of this chapter:} \emph{I can tell how a thief broke in and was caught: to break open, a bolt, to snatch, to catch in the act and to arrest.}

Complete the last lesson of chapter 304: I can tell how a thief broke in and was caught: to break open, a bolt, to snatch, to catch in the act and to arrest.
\end{chapteropening}

\section[effringō, effringere]{effringō, effringere — to break open}

\label{lesson:LA-C304-effringo}



\begin{quote}
Before the new one: say the Latin for a robber, then the Latin for loot.
\end{quote}

\subsection*{You'll want to know: effringō, effringere}
\textbf{effringō, effringere} — \textquotedblleft{}to break open\textquotedblright{}.

\begin{quote}
Fūr portam effrēgit et in domum intrāvit.
\end{quote}

\begin{quote}
\emph{The thief broke the door open and went into the house.}
\end{quote}

\emph{effringō, effringere, effrēgī, effrāctum}: \emph{ex} + \emph{frangō} (taught). Plautus has doors smashed in this way.

\begin{cousinweb}
\textbf{effringō} is built on \emph{frangō}, \textquotedblleft{}to break\textquotedblright{}, with \emph{ex-}, \textquotedblleft{}out\textquotedblright{} (Lewis \& Short); the vowel changes inside the compound.
\end{cousinweb}

\begin{grammarlens}[title={What you've built}]
One more word for telling, explaining and describing at B1.
\end{grammarlens}

\subsection*{Your turn}
\begin{itemize}
\raggedright
  \item \emph{Say it:} \emph{effringō, effringere}
  \item \emph{Say it:} \emph{effringō, effringere}, once more
  \item \emph{Say it:} a sentence of your own with \emph{effringō, effringere}
  \item \emph{Recall:} say the Latin for a robber, then the Latin for loot, then say \emph{effringō, effringere} again
\end{itemize}

\begin{tcolorbox}[breakable,colback=teal!4,colframe=teal!35!black,title={Before you move on}]
What does effringō, effringere mean? (\textquotedblleft{}To break open\textquotedblright{}.) Say it once more.
\end{tcolorbox}
\section[claustrum, claustrī]{claustrum, claustrī — a bolt, a bar; a barrier}

\label{lesson:LA-C304-claustrum}



\begin{quote}
Before the new one: say the Latin for loot, then the Latin for to break open.
\end{quote}

\subsection*{You'll want to know: claustrum, claustrī}
\textbf{claustrum, claustrī} — \textquotedblleft{}a bolt, a bar; a barrier\textquotedblright{}.

\begin{quote}
Fūr claustra portae nocte remōvit.
\end{quote}

\begin{quote}
\emph{At night the thief removed the bolts of the door.}
\end{quote}

Neuter, from \emph{claudō}, 'to shut' (taught); most often plural, \emph{claustra}, the bars of a door or gate. English \emph{cloister}.

\begin{cousinweb}
\textbf{claustrum} is built on \emph{claudō}, \textquotedblleft{}to shut\textquotedblright{} (EtymDB), with a tool ending: the thing that shuts. English \emph{cloister} comes from it through French, and \emph{claustrophobia} is built on it.
\end{cousinweb}

\begin{grammarlens}[title={What you've built}]
One more word for telling, explaining and describing at B1.
\end{grammarlens}

\subsection*{Your turn}
\begin{itemize}
\raggedright
  \item \emph{Say it:} \emph{claustrum, claustrī}
  \item \emph{Say it:} \emph{claustrum, claustrī}, once more
  \item \emph{Say it:} a sentence of your own with \emph{claustrum, claustrī}
  \item \emph{Recall:} say the Latin for loot, then the Latin for to break open, then say \emph{claustrum, claustrī} again
\end{itemize}

\begin{tcolorbox}[breakable,colback=teal!4,colframe=teal!35!black,title={Before you move on}]
What does claustrum, claustrī mean? (\textquotedblleft{}A bolt, a bar; a barrier\textquotedblright{}.) Say it once more.
\end{tcolorbox}
\section[rapiō, rapere]{rapiō, rapere — to snatch, to carry off}

\label{lesson:LA-C304-rapio}



\begin{quote}
Before the new one: say the Latin for to break open, then the Latin for a bolt.
\end{quote}

\subsection*{You'll want to know: rapiō, rapere}
\textbf{rapiō, rapere} — \textquotedblleft{}to snatch, to carry off\textquotedblright{}.

\begin{quote}
Fūr crumēnam meam rapuit et fūgit.
\end{quote}

\begin{quote}
\emph{The thief snatched my purse and ran away.}
\end{quote}

\emph{rapiō, rapere, rapuī, raptum}. English \emph{rapid} and \emph{rapture} come from it.

\begin{cousinweb}
\textbf{rapiō} is inherited from Proto-Italic \emph{*rapiō} (EtymDB). English \emph{rapid} comes from \emph{rapidus}, \textquotedblleft{}snatching, swift\textquotedblright{}, and \emph{rapture} from its participle.
\end{cousinweb}

\begin{grammarlens}[title={What you've built}]
One more word for telling, explaining and describing at B1.
\end{grammarlens}

\subsection*{Your turn}
\begin{itemize}
\raggedright
  \item \emph{Say it:} \emph{rapiō, rapere}
  \item \emph{Say it:} \emph{rapiō, rapere}, once more
  \item \emph{Say it:} a sentence of your own with \emph{rapiō, rapere}
  \item \emph{Recall:} say the Latin for to break open, then the Latin for a bolt, then say \emph{rapiō, rapere} again
\end{itemize}

\begin{tcolorbox}[breakable,colback=teal!4,colframe=teal!35!black,title={Before you move on}]
What does rapiō, rapere mean? (\textquotedblleft{}To snatch, to carry off\textquotedblright{}.) Say it once more.
\end{tcolorbox}
\section[dēprehendō, dēprehendere]{dēprehendō, dēprehendere — to catch, to catch in the act}

\label{lesson:LA-C304-deprehendo}



\begin{quote}
Before the new one: say the Latin for a bolt, then the Latin for to snatch.
\end{quote}

\subsection*{You'll want to know: dēprehendō, dēprehendere}
\textbf{dēprehendō, dēprehendere} — \textquotedblleft{}to catch, to catch in the act\textquotedblright{}.

\begin{quote}
Servus fūrem in hortō dēprehendit.
\end{quote}

\begin{quote}
\emph{A slave caught the thief in the garden.}
\end{quote}

\emph{dēprehendō, dēprehendere, dēprehendī, dēprehēnsum}, also \emph{dēprendō}. \emph{In fūrtō dēprehendere}: to catch someone stealing.

\begin{cousinweb}
\textbf{dēprehendō} is built on \emph{prehendō}, \textquotedblleft{}to seize\textquotedblright{}, with \emph{dē-} (Lewis \& Short): to snatch away, to catch something on the move.
\end{cousinweb}

\begin{grammarlens}[title={What you've built}]
One more word for telling, explaining and describing at B1.
\end{grammarlens}

\subsection*{Your turn}
\begin{itemize}
\raggedright
  \item \emph{Say it:} \emph{dēprehendō, dēprehendere}
  \item \emph{Say it:} \emph{dēprehendō, dēprehendere}, once more
  \item \emph{Say it:} a sentence of your own with \emph{dēprehendō, dēprehendere}
  \item \emph{Recall:} say the Latin for a bolt, then the Latin for to snatch, then say \emph{dēprehendō, dēprehendere} again
\end{itemize}

\begin{tcolorbox}[breakable,colback=teal!4,colframe=teal!35!black,title={Before you move on}]
What does dēprehendō, dēprehendere mean? (\textquotedblleft{}To catch, to catch in the act\textquotedblright{}.) Say it once more.
\end{tcolorbox}
\section[comprehendō, comprehendere]{comprehendō, comprehendere — to seize, to arrest; to grasp}

\label{lesson:LA-C304-comprehendo}



\begin{quote}
Before the new one: say the Latin for to snatch, then the Latin for to catch.
\end{quote}

\subsection*{You'll want to know: comprehendō, comprehendere}
\textbf{comprehendō, comprehendere} — \textquotedblleft{}to seize, to arrest; to grasp\textquotedblright{}.

\begin{quote}
Mīlitēs latrōnem comprehendērunt et in carcerem dūxērunt.
\end{quote}

\begin{quote}
\emph{The soldiers arrested the bandit and took him to prison.}
\end{quote}

\emph{comprehendō, comprehendere, comprehendī, comprehēnsum}. To seize a person, and also to grasp an idea. English \emph{comprehend}.

\begin{cousinweb}
\textbf{comprehendō} is built on \emph{prehendō}, \textquotedblleft{}to seize\textquotedblright{}, with \emph{con-}, \textquotedblleft{}together, wholly\textquotedblright{} (EtymDB): to take hold of on all sides. English \emph{comprehend} comes from it.
\end{cousinweb}

\begin{grammarlens}[title={What you've built}]
One more word for telling, explaining and describing at B1.
\end{grammarlens}

\subsection*{Your turn}
\begin{itemize}
\raggedright
  \item \emph{Say it:} \emph{comprehendō, comprehendere}
  \item \emph{Say it:} \emph{comprehendō, comprehendere}, once more
  \item \emph{Say it:} a sentence of your own with \emph{comprehendō, comprehendere}
  \item \emph{Recall:} say the Latin for to snatch, then the Latin for to catch, then say \emph{comprehendō, comprehendere} again
\end{itemize}

\begin{tcolorbox}[breakable,colback=teal!4,colframe=teal!35!black,title={Before you move on}]
What does comprehendō, comprehendere mean? (\textquotedblleft{}To seize, to arrest; to grasp\textquotedblright{}.) Say it once more.
\end{tcolorbox}
